% 柯伊伯带（综述）
% license CCBYSA3
% type Wiki

本文根据 CC-BY-SA 协议转载翻译自维基百科\href{https://en.wikipedia.org/wiki/Kuiper_belt}{相关文章}。

柯伊伯带（/ˈkaɪpər/ ⓘ）$^{[2]}$ 是位于太阳系外部的一种环星盘结构，其范围从海王星轨道处约 30 个天文单位（AU）一直延伸到距离太阳约 50 AU 的区域。$^{[3]}$ 它在结构上类似于小行星带，但规模要大得多——宽度约为其 20 倍，总质量约为其 20–200 倍。$^{[4][5]}$ 与小行星带相似，柯伊伯带主要由太阳系形成早期遗留下来的小天体或残骸组成。不同之处在于，许多小行星主要由岩石和金属构成，而大多数柯伊伯带天体则主要由冻结的挥发性物质（通常称为“冰”）构成，例如甲烷、氨和水冰。柯伊伯带是天文学界普遍认可的多颗矮行星的家园，包括 Orcus、冥王星（Pluto）$^{[6]}$、妊神星（Haumea）$^{[7]}$、夸奥瓦（Quaoar）以及鸟神星（Makemake）$^{[8]}$。太阳系中的一些卫星，例如海王星的海卫一（Triton）和土星的菲比（Phoebe），也可能起源于这一地区。$^{[9][10]}$

柯伊伯带以荷兰天文学家杰拉德·柯伊伯（Gerard Kuiper）的名字命名，他在 1951 年提出了这种带状结构存在的一种设想。$^{[11]}$ 在他之前和之后，也有研究者提出过类似的假说，例如 20 世纪 30 年代的肯尼思·埃奇沃斯（Kenneth Edgeworth）。$^{[12]}$ 对该结构最为直接的预测来自天文学家胡利奥·安赫尔·费尔南德斯（Julio Ángel Fernández），他在 1980 年发表论文，提出在海王星之外存在一条彗星带$^{[13][14]}$，并认为它可能是短周期彗星的来源。$^{[15][16]}$

1992 年，小行星 15760 Albion 被发现，这是继 1930 年发现冥王星、1978 年发现卡戎之后确认的首个柯伊伯带天体（KBO）。$^{[17]}$ 自其发现以来，已知的柯伊伯带天体数量已增加到数千个，并且人们认为直径超过 100 km（62 英里）的 KBO 可能多达 100000 个以上。$^{[18]}$ 起初，柯伊伯带被认为是周期彗星（轨道周期小于 200 年）的主要储存库。然而，自 20 世纪 90 年代中期以来的研究表明，柯伊伯带在动力学上是稳定的，而彗星真正的起源地是散射盘——这是一个由海王星在约 45 亿年前向外迁移所形成的动力学活跃区域。$^{[19]}$ 散射盘天体（例如阋神星 Eris）具有极高的轨道偏心率，其轨道最远可延伸至距离太阳约 100 AU 的位置。$^{[a]}$

柯伊伯带不同于假想中的奥尔特云，后者被认为距离太阳远得多（约为其一千倍），且整体上近似球形。柯伊伯带内的天体，与散射盘成员以及任何潜在的希尔斯云或奥尔特云天体，统称为海王星外天体（TNO）。$^{[22]}$ 冥王星是目前已知柯伊伯带中最大、质量最大的成员，同时也是已知 TNO 中体积最大、质量第二大的天体，仅次于散射盘中的阋神星。$^{[a]}$ 冥王星最初被视为一颗行星，但由于其属于柯伊伯带成员，这一事实促使其在 2006 年被重新归类为矮行星。冥王星在成分上与许多其他柯伊伯带天体相似，其轨道周期也具有一类 KBO 的典型特征，这类天体被称为“冥族天体”（plutinos），它们与海王星处于相同的 2:3 轨道共振状态。

柯伊伯带与海王星有时也被视为太阳系范围的标志之一，其他替代界限还包括日球层顶（heliopause），以及太阳引力影响与其他恒星引力影响相当的距离（估计约在 50000–125000 AU 之间）。$^{[23]}$
\subsection{历史}
\begin{figure}[ht]
\centering
\includegraphics[width=6cm]{./figures/c3c318a9ff09ac52.png}
\caption{冥王星与卡戎} \label{fig_Kuiper_1}
\end{figure}
1930 年冥王星被发现之后，许多人开始推测它可能并非孤立存在。如今被称为柯伊伯带的这一区域，在随后数十年间以多种不同形式被提出为假想结构。直到 1992 年，人们才首次获得其存在的直接证据。由于在此之前关于柯伊伯带性质的设想在数量和形式上都十分多样，至今仍难以确定究竟应将最早提出这一概念的功劳归于谁。$^{[24]:106}$
\subsubsection{假说}
最早提出海王星外天体族群存在设想的天文学家是弗雷德里克·C·伦纳德。在 1930 年克莱德·汤博发现冥王星之后不久，伦纳德便思考，冥王星是否“很可能只是一系列超海王星天体中被发现的第一个，其余成员仍有待发现，但终将被探测到”。$^{[25]}$ 同年，天文学家阿尔明·O·洛伊施纳也提出，冥王星“可能是许多尚未发现的长周期行星状天体之一”。$^{[26]}$
\begin{figure}[ht]
\centering
\includegraphics[width=6cm]{./figures/acd1d13b7d423c69.png}
\caption{以其命名柯伊伯带的天文学家杰拉德·柯伊伯} \label{fig_Kuiper_2}
\end{figure}
1943 年，肯尼思·埃奇沃斯（Kenneth Edgeworth）在《英国天文学会期刊》中提出假说，认为在海王星之外的区域，原始太阳星云中的物质分布过于稀疏，无法凝聚形成行星，而是凝结成数量众多的小天体。由此他得出结论认为，“太阳系的外部区域，即行星轨道之外，被大量相对较小的天体所占据”$^{[27]:xii}$，并且这些天体中的个别成员会“不时偏离自身的活动范围，作为偶然的访客进入内太阳系”$^{[27]:2}$，从而表现为彗星。

1951 年，杰拉德·柯伊伯（Gerard Kuiper）在《天体物理学：专题研讨会》中的一篇论文里，对太阳系早期演化过程中形成的类似盘状结构进行了推测，并认为该盘由“原始团块的残余组成，其中许多成员已经流失，成为游离的小行星，类似于疏散的银河开放星团解体为恒星的过程”$^{[11]}$。在另一篇基于他 1950 年所作讲座、同样题为《论太阳系的起源》的论文中，柯伊伯描述了“太阳星云最外侧区域，从 38 到 50 个天文单位（即原始海王星轨道之外）”，并指出在那里“凝结产物（如 H(_2)O、NH(_3)、CH(_4) 等冰）必然形成，这些冰片会缓慢聚集，形成更大的集合体，其尺度估计可达 1 km 甚至更大”。他进一步指出，“这些凝结体在尺寸、数量和组成上似乎可以解释彗星的性质”。根据柯伊伯的观点，“冥王星在 30 到 50 个天文单位的整个区域内运行，应当被认为是启动彗星向整个太阳系散射的原因”$^{[28]}$。柯伊伯当时基于一个在其时代较为普遍的假设，即冥王星的质量远大于我们今天所知道的数值，因此能够将这些天体散射至奥尔特云甚至抛射出太阳系；如果这一假设成立，那么今天将不会存在柯伊伯带$^{[29]}$。

在随后的几十年中，这一假说以多种形式不断演化。1962 年，物理学家阿拉斯泰尔·G·W·卡梅伦（Alastair G. W. Cameron）提出，在太阳系边缘存在“一团极其巨大的小尺度物质”$^{[27]:14}$。1964 年，弗雷德·惠普尔（Fred Whipple）——因提出著名的“脏雪球”彗星结构模型而广为人知——认为存在一条“彗星带”，其质量可能大到足以解释此前促使人们寻找 X 行星的天王星轨道异常，或至少足以显著影响已知彗星的轨道$^{[30]}$。然而，后续观测排除了这一假说$^{[27]:14}$。

1977 年，查尔斯·科瓦尔（Charles Kowal）发现了 2060 Chiron，这是一颗具有冰质成分、轨道位于土星与天王星之间的天体。他使用了闪烁比较仪——正是这一仪器在近 50 年前帮助克莱德·汤博发现了冥王星$^{[31]}$。1992 年，又发现了另一颗具有类似轨道的天体 5145 Pholus$^{[32]}$。如今，人们已经确认在木星与海王星之间存在着一整类类似彗星的天体群体，被称为半人马天体。半人马天体的轨道是不稳定的，其动力学寿命仅有数百万年$^{[33]}$。自 1977 年发现 Chiron 以来，天文学家便推测，这些半人马天体必然持续从某个外部储库中得到补充$^{[27]:38}$。

关于柯伊伯带存在的进一步证据后来来自对彗星的研究。人们早已认识到，彗星具有有限的寿命。当它们接近太阳时，太阳辐射的热量会使其表面的挥发性物质升华并逸散到太空中，从而逐渐消散。为了使彗星在整个太阳系年龄尺度上持续可见，它们必须不断得到补充$^{[34]}$。一种被提出的补给区域是假想中的奥尔特云，这可能是一个球状的彗星云团，延伸至距太阳 50 000 AU 之外，由荷兰天文学家扬·奥尔特（Jan Oort）于 1950 年首次提出$^{[35]}$。奥尔特云被认为是长周期彗星的起源地，这类彗星（如海尔–波普彗星）的轨道周期可长达数千年$^{[24]:105}$。
\begin{figure}[ht]
\centering
\includegraphics[width=6cm]{./figures/06da66ebf1035986.png}
\caption{1980 年，天文学家胡利奥·费尔南德斯（Julio Fernández）预测了这样一条天体带的存在。有人指出，由于在费尔南德斯论文的开篇句中同时出现了“Kuiper”和“comet belt（彗星带）”这两个词，这一假想中的区域因此被称为“柯伊伯带”$^{[36]}$。} \label{fig_Kuiper_3}
\end{figure}
还存在另一类彗星族群，被称为短周期彗星或周期彗星，其成员的轨道周期小于 200 年，例如哈雷彗星。到 20 世纪 70 年代，人们发现短周期彗星的发现速率越来越难以与其完全源自奥尔特云这一假设相一致。$^{[27]:39}$ 对于一个奥尔特云天体而言，若要成为短周期彗星，首先必须被巨行星俘获。1980 年，乌拉圭天文学家胡利奥·费尔南德斯（Julio Fernández）在《英国皇家天文学会月报》上发表论文指出：每向内太阳系输送 1 颗短周期彗星，就必须有约 600 颗天体被抛射到星际空间中。他据此推测，为了解释观测到的彗星数量，需要在 35–50 AU 之间存在一条彗星带。$^{[37]}$在费尔南德斯工作的基础上，1988 年，加拿大的马丁·邓肯（Martin Duncan）、汤姆·奎因（Tom Quinn）和斯科特·特雷梅恩（Scott Tremaine）团队进行了一系列计算机模拟，以检验所有已观测到的彗星是否都可能来自奥尔特云。他们发现，奥尔特云无法解释全部短周期彗星，尤其是因为短周期彗星往往聚集在太阳系的黄道平面附近，而奥尔特云彗星则可能从天空中的任意方向进入内太阳系。若在模型中加入费尔南德斯所设想的这条“带状结构”，模拟结果便能够与观测相符合。$^{[38]}$ 据称，正是因为费尔南德斯论文的开篇句中同时出现了“Kuiper”和“comet belt”这两个词，特雷梅恩将这一假想区域命名为“柯伊伯带”。$^{[27]:191}$
\subsubsection{发现}
\begin{figure}[ht]
\centering
\includegraphics[width=8cm]{./figures/a001aa6351c120b0.png}
\caption{位于冒纳凯阿山顶的望远镜群。柯伊伯带正是通过 UH88 望远镜发现的，它是从左数第四台。} \label{fig_Kuiper_4}
\end{figure}
1987 年，当时在麻省理工学院（MIT）任职的天文学家戴维·朱伊特（David Jewitt）对“外太阳系显得明显空旷”这一现象愈发感到困惑 $^{[17]}$。他鼓励当时仍是研究生的简·露（Jane Luu）协助他开展寻找冥王星轨道之外天体的工作，因为正如他对她所说的那样：“如果我们不去做，就不会有人去做。”$^{[27]:50}$。朱伊特和露利用位于美国亚利桑那州的基特峰国家天文台以及智利的托洛洛山美洲际天文台的望远镜，采用与克莱德·汤博和查尔斯·科瓦尔当年几乎相同的方法开展搜索，即使用闪烁比较仪进行比对 $^{[27]:50}$。最初，检查每一对底片大约需要八个小时 $^{[27]:51}$，但随着电子电荷耦合器件（CCD）的出现，这一过程显著加快。尽管 CCD 的视场较窄，但其集光效率远高于传统摄影底片（可保留约 90\% 的入射光，而摄影底片仅约 10\%），并且允许在计算机屏幕上以电子方式完成“闪烁”比较。如今，CCD 已成为大多数天文探测器的基础 $^{[27]:52,54,56}$。
1988 年，朱伊特转至夏威夷大学天文学研究所工作，露随后也加入他，在毛纳基山使用夏威夷大学的 2.24 米望远镜开展观测 $^{[27]:57,62}$。随着 CCD 视场最终扩大到 1024×1024 像素，搜索效率得以大幅提升 $^{[27]:65}$。最终，在历时五年的搜索之后，朱伊特和露于 1992 年 8 月 30 日宣布“发现候选柯伊伯带天体 1992 QB1” $^{[17]}$。该天体后来被命名为 15760 Albion。六个月后，他们又在该区域发现了第二个天体 (181708) 1993 FW $^{[39]}$。截至 2018 年，已发现的柯伊伯带天体数量超过 2000 个 $^{[40]}$。

在发现 1992 QB1（2018 年命名为 15760 Albion）之后的二十年间（1992—2012 年），在这一带区共发现了一千多个天体，显示出除冥王星和 Albion 之外还存在着一个极为庞大的天体带 $^{[41]}$。即便在 2010 年代，人们对柯伊伯带天体的整体范围和性质仍知之甚少 $^{[41]}$。最终，无人探测器“新视野号”（New Horizons）完成了首次柯伊伯带天体的近距离飞掠观测，先后对冥王星系统（2015 年）以及 Arrokoth（2019 年）进行了更为细致的观测 $^{[42]}$。

自跨海王星区域首次被系统绘制以来的研究表明，如今所谓的柯伊伯带并非短周期彗星的真正起源地，这些彗星实际上源自一个相关联的天体族群，即散射盘。散射盘形成于海王星向外迁移进入原始柯伊伯带的过程中，当时该区域距离太阳更近。海王星的迁移在其身后留下了两类天体：一类是动力学上稳定、其轨道不再受海王星影响的天体（即严格意义上的柯伊伯带），另一类是近日点足够接近、仍会在海王星绕日运行过程中受到其扰动的天体（即散射盘）。由于散射盘在动力学上更为活跃，而柯伊伯带相对稳定，散射盘如今被认为是周期彗星最可能的起源地 $^{[19]}$。
\subsubsection{名称}
天文学家有时使用“埃奇沃思—柯伊伯带”（Edgeworth–Kuiper belt）这一替代名称以表彰埃奇沃思的贡献，柯伊伯带天体（KBO）有时也被称为 EKOs。布赖恩·G·马斯登（Brian G. Marsden）认为二者都不应被视为真正的提出者：“无论是埃奇沃思还是柯伊伯，都没有描述过任何与我们今天所看到的情形真正相似的东西，真正做到这一点的是弗雷德·惠普尔。”$^{[27]:199}$。戴维·朱伊特则评论道：“如果一定要说的话……最接近成功预测柯伊伯带的人其实是费尔南德斯。”$^{[29]}$

柯伊伯带天体有时也被称为“kuiperoids”，这一名称由克莱德·汤博提出 $^{[43]}$。若干科学团体建议使用“跨海王星天体”（trans-Neptunian object，TNO）这一术语来指代该带区内的天体，因为该术语争议较少。不过，两者并非严格同义：TNO 指的是所有绕太阳运行且轨道位于海王星轨道之外的天体，并不仅限于柯伊伯带天体 $^{[44]}$。
\subsection{结构}
在其最大范围内（包括外围区域，但不包括散射盘），柯伊伯带大致从距太阳约 30–55 AU 的区域延伸。通常认为，柯伊伯带的主体部分从 39.5 AU 处的 2:3 平均运动共振（见下文）一直延伸到约 48 AU 处的 1:2 共振位置 $^{[45]}$。柯伊伯带在空间中相当“厚”，其主要集中区可延伸到黄道面上下约 10°，而更为稀疏的天体分布则可延伸到数倍于此的角距离。整体而言，它看起来更像一个环面或甜甜圈，而非一条狭长的带状结构 $^{[46]}$。其平均位置相对于黄道面的倾角约为 1.86° $^{[47]}$。

海王星的存在通过轨道共振对柯伊伯带的结构产生了深远影响。在与太阳系年龄相当的时间尺度上，海王星的引力会使位于某些特定区域内的天体轨道变得不稳定，并将它们要么送入内太阳系，要么抛入散射盘，甚至逐出星际空间。这一过程使柯伊伯带在当前结构中出现了明显的空隙，类似于小行星带中的柯克伍德空隙。例如，在 40–42 AU 之间的区域，天体无法长期保持稳定轨道，因此任何在该区域内观测到的天体都必然是相对近期才迁移到那里的 $^{[48]}$。
\begin{figure}[ht]
\centering
\includegraphics[width=10cm]{./figures/78989492521f7356.png}
\caption{柯伊伯带的半长轴与偏心率结构示意图。冥族天体（plutinos）以橙色表示，其他处于共振状态的天体以红色表示。非共振的“热”经典柯伊伯带天体（hot cubewanos）以天蓝色表示，而“冷”经典柯伊伯带天体（cold cubewanos）以深蓝色表示。妊神星（Haumea）碰撞家族以蓝灰色表示。图中标注了若干重要天体。} \label{fig_Kuiper_5}
\end{figure}
\subsubsection{经典带}
在与海王星形成 2:3 与 1:2 平均运动共振之间、约位于 42–48 AU 的区域内，与海王星的引力相互作用发生在较长的时间尺度上，因此天体可以在轨道几乎不被改变的情况下长期存在。该区域被称为经典柯伊伯带，其成员约占目前已观测到的柯伊伯带天体（KBO）的三分之二$^{[49][50]}$。由于现代意义上发现的第一个柯伊伯带天体（Albion，长期被称为 1992 QB1）被视为这一类的原型，经典 KBO 常被称为 cubewanos（“Q-B-1-os”）$^{[51][52]}$。国际天文学联合会（IAU）制定的命名规范要求，经典 KBO 应以与创世相关的神话人物命名$^{[53]}$。

经典柯伊伯带似乎由两个不同的天体族群组成。第一类被称为“动力学冷”族群，其轨道与行星相似：轨道几乎为圆形，偏心率小于 0.1，轨道倾角也较低，最高约为 10°（即接近太阳系的黄道面，而非显著倾斜）。冷族群中还包含一组被称为“核心（kernel）”的天体，其半长轴集中在 44–44.5 AU$^{[54]}$。第二类是“动力学热”族群，其轨道相对于黄道面的倾角可高达 30°。这两个族群的命名并非源于温度的显著差异，而是借用了气体粒子的类比：粒子在被“加热”时，其相对速度会增大$^{[55]}$。

这两个族群不仅在轨道性质上不同，冷族群在颜色和反照率方面也存在差异：它们通常颜色更红、亮度更高，具有更高比例的双星系统$^{[56]}$，其尺寸分布也不同$^{[57]}$，且几乎不包含非常大的天体$^{[58]}$。动力学冷族群的总质量大约只有热族群的三十分之一$^{[57]}$。颜色差异可能反映了成分上的不同，这表明它们形成于不同的区域。动力学热族群被认为形成于海王星最初的轨道附近，并在巨行星迁移过程中被散射到外侧$^{[4][59]}$；而动力学冷族群则被认为大致形成于其当前位置，因为较为松散的双星系统在与海王星的近距离遭遇中很难幸存$^{[60]}$。尽管“尼斯模型”似乎能够至少部分解释这种成分差异，但也有人提出，颜色差异可能反映的是天体表面演化过程的不同$^{[61]}$。
\subsubsection{共振}
\begin{figure}[ht]
\centering
\includegraphics[width=8cm]{./figures/17c17f1d2942df42.png}
\caption{跨海王星天体的分布图，突出显示了轨道共振关系。其中，冥族天体（Plutinos）以橙色表示，其他处于共振状态的天体以红色表示；非共振的 cubewanos 以及散射盘天体分别以蓝色和灰色表示。} \label{fig_Kuiper_6}
\end{figure}
\begin{figure}[ht]
\centering
\includegraphics[width=6cm]{./figures/22b1e1c1bc4ee3d5.png}
\caption{} \label{fig_Kuiper_7}
\end{figure}

\subsection{另见}
\begin{itemize}
\item 小行星带
\item 可能的矮行星列表
\item 海王星外天体列表
\item 第九行星
\item 海王星外天体光谱类型
\end{itemize}
\subsection{注释}
a.关于术语“散射盘”和“柯伊伯带”的使用，文献中存在不一致之处。对一些人而言，这两个概念是截然不同的天体群；而另一些人则认为，散射盘是柯伊伯带的一部分。甚至有作者在同一篇出版物中会在这两种用法之间交替使用。由于国际天文学联合会下属的小行星中心——负责编目太阳系内小行星的机构——做出了这一区分，维基百科关于跨海王星区域的文章在编辑时也采用了同样的区分。在维基百科上，已知质量最大的海王星外天体厄里斯并不属于柯伊伯带，这使得冥王星成为质量最大的柯伊伯带天体。
\subsection{参考文献}
\begin{enumerate}
\item 来源：小行星中心, www.cfeps.net等。
\item “柯伊伯带”.Dictionary.com 精编版（在线）。无日期。；Merriam-Webster.com 词典。梅里亚姆-韦伯斯特。“发音提示：快速参考”Merriam-Webster.com梅里亚姆-韦伯斯特已归档自原始网页于2007年6月1日存档。 检索于2025年4月3日
\item 斯特恩，艾伦；科尔韦尔，乔舒亚·E.（1997）。“原始埃奇沃斯-柯伊伯带中的碰撞侵蚀与30至50天文单位柯伊伯空隙的形成”. 天体物理学报. 490 (2): 879–882. 文献标识码:1997ApJ...490..879S. doi:10.1086/304912.
\item 德尔桑蒂，奥黛丽 & 朱维特，大卫（2006）。行星之外的太阳系（PDF）。夏威夷大学天文学研究所。文献标识码:2006ssu..book..267D。于2007年9月25日从原文存档。 检索于2007年3月9日。
\item 克拉斯金斯基，G. A.; 彼捷娃，E. V.；瓦西里耶夫，M. V.；亚古季娜，E. I. （2002年7月）。“小行星带中的暗物质”。国际行星科学杂志. 158 (1): 98–105. 天文学数据库标识符:2002Icar..158...98K. DOI:10.1006/icar.2002.6837.
\item 克里斯滕森，拉斯·林德伯格。“国际天文学联合会2006年大会：国际天文学联合会决议投票结果”。国际天文学联合会。于2014年4月29日存档自原始网页。检索日期：2021年5月25日。
\item 克里斯滕森，拉斯·林德伯格（2008年9月17日）。“国际天文学联合会命名第五颗矮行星为妊神星”。国际天文学联合会。于2014年4月25日存档自原文。检索日期：2021年5月25日。
\item 克里斯滕森，拉斯·林德伯格（2008年7月19日）。“第四颗矮行星被命名为马凯马克”。国际天文学联合会。于2019年6月16日存档自原文。检索日期：2021年5月25日。
\item 约翰逊，托伦斯·V.；以及卢宁，乔纳森·I.土星的卫星菲比作为来自太阳系外层的捕获天体, 自然，第435卷，第69–71页。
\item 克雷格·B·阿格诺尔与道格拉斯·P·汉密尔顿（2006）。“海王星在双行星引力遭遇中捕获其卫星海卫一”（PDF）。自然441（7090）：192–194。天文学编码：2006Natur.441..192ADOI10.1038/nature04792PMID16688170S2CID4420518。于2007年6月21日从原始网页存档。检索日期：2006年6月20日
\item 开普勒，G. P.（1951）。“关于太阳系的起源”。载于海内克，J. A.（编）。天体物理学：专题研讨会。美国纽约市，纽约州：麦格劳-希尔出版社。第357–424.
\item “柯伊伯带：事实——美国国家航空航天局科学”。2017年11月14日。于2024年3月12日存档自原始网页。检索日期：2024年2月18日。
\item J. A. 费尔南德斯（1980）。“关于海王星外彗星带的存在”. 马德里国家天文台. 192 (3): 481–491. Bibcode:1980MNRAS.192..481F. doi:10.1093/mnras/192.3.481.
\item 莫尔比戴利，A.；托马斯，F.；穆恩斯，M.（1995年12月1日）。“柯伊伯带的共振结构与前五个海王星外天体的动力学”. 《伊卡洛斯》. 118 (2): 322–340. 文献标识码:1995Icar..118..322M. DOI:10.1006/icar.1995.1194. ISSN 0019-1035.
\item gunjan.sogani（2022年9月10日）。“柯伊伯带及其成员的发现”。Wondrium每日。于2023年8月1日从原文存档。检索日期：2023年8月1日
\item “胡里奥·A·费尔南德斯”.美国国家科学院官网于2023年8月1日存档自原始网页。检索日期：2023年8月1日
\item 朱维特，大卫；刘佳（1993）。“柯伊伯带候选天体1992 QB1的发现”。自然.362(6422)：730–732。Bibcode：1993Natur.362..730Jdoi10.1038/362730a0S2CID4359389
\item “首席研究员的观点”.新视野。2012年8月24日。于2014年11月13日从原始网页存档.
\item 莱维森，哈罗德·F.；多尼斯，卢克（2007）。“彗星种群与彗星动力学”。载于露西·安·亚当斯·麦克法登；保罗·罗伯特·魏斯曼；托伦斯·V·约翰逊（编）。《太阳系百科全书》（第2版）。阿姆斯特丹；波士顿：学术出版社。第575–588页.ISBN978-0-12-088589-3
\item 魏斯曼与约翰逊，2007年，《太阳系百科全书》，脚注第584页。
\item “半人马星与散盘天体列表”.国际天文学联合会：小行星中心. 哈佛-史密森天体物理中心中央天文电报局。2011年1月3日。已归档于2017年6月29日自原始网页存档。检索于2011年1月3日
\item 热拉尔·福尔（2004）。“截至2004年5月20日的小行星系统描述”。于2007年5月29日从原文存档。检索日期：2007年6月1日。
\item “太阳系的边缘在哪里？”.戈达德媒体工作室. 美国国家航空航天局戈达德太空飞行中心。2017年9月5日。已存档自原始网页于2021年12月16日存档于2019年9月22日
\item 兰德尔，丽莎（2015）。暗物质与恐龙。纽约：埃科/哈珀柯林斯出版社。ISBN 978-0-06-232847-2.
\item “‘柯伊伯带’这一术语有何不妥之处？（或者，为何要以一位不相信其存在的人的名字来命名这一天体？）”。国际彗星季刊于2019年10月8日存档自原始网页。检索日期：2010年10月24日
\item 戴维斯，约翰·K.；麦克法兰，J.；贝利，马克·E.；马斯登，布赖恩·G.；叶伟文（2008）。“关于跨海王星区域的早期思想发展”（PDF）。载于M. 安东内塔·巴拉奇；赫尔曼·博恩哈特；戴尔·克鲁伊昌克；亚历山德罗·莫尔比德利（编）。
\item 戴维斯，约翰·K.（2001）。超越冥王星：探索太阳系的外缘。剑桥大学出版社。
\item 开普勒，杰拉德（1951）。“论太阳系的起源”. 美国国家科学院院刊. 37 (1): 233. doi:10.1073/pnas.37.4.233. PMC 1063291. PMID 16588984.
\item 大卫·朱维特。“为什么叫‘柯伊伯’带？”。夏威夷大学于2019年2月12日从原始网页存档。检索日期：2007年6月14日
\item 饶，M. M.（1964）。“向量测度的分解”（PDF）。美国国家科学院院刊51（5）：771–774。文献标识码：1964PNAS...51..771RDOI10.1073/pnas.51.5.771PMC300359PMID16591174已归档自原始网页于2016年6月3日。检索于2007年6月20日
\item C. T. 科瓦尔；W. 利勒；B. G. 马尔斯登（1977）。“/2060/ 凯龙的发现与轨道”。载于：太阳系动力学；研讨会论文集. 81。哈雷天文台，哈佛—史密森天体物理中心：245页。文献编码:1979IAUS...81..245K.
\item J. V. 斯科蒂；D. L. 拉比诺维茨；C. S. 肖梅克；E. M. 肖梅克；D. H. 利维；T. M. 金；E. F. 海林；J. 阿卢；K. 劳伦斯；R. H. 麦克诺特；L. 弗雷德里克；D. 托伦；B. E. A. 穆勒（1992）。“1992 AD”。国际天文学联合会通告 (5434): 1. 文献标识码:1992IAUC.5434....1S.
\item 霍纳，J.；埃文斯，N. W.；贝利，M. E.（2004）。“半人马星群的模拟Ⅰ：总体统计”. 皇家天文学会月报. 354 (3): 798–810. arXiv:astro-ph/0407400. 文献编码:2004MNRAS.354..798H. DOI:10.1111/j.1365-2966.2004.08240.x. S2CID 16002759.
\item 大卫·朱维特（2002）。“从柯伊伯带天体到彗星核：缺失的远红外观测物质”. 《天文期刊》. 123 (2): 1039–1049. 文献标识码:2002AJ....123.1039J. DOI:10.1086/338692. 跨学科索引 122240711.
\item 奥尔特，J. H.（1950）。“环绕太阳系的彗星云的结构及其起源假说”。荷兰天文研究所公报. 11: 91. 文献标识码:1950BAN....11...91O.
\item “柯伊伯带”。九大行星。2019年10月8日。于2021年5月16日存档自原始网页。检索日期：2020年8月16日
\item J. A. 费尔南德斯（1980）。“关于海王星外侧存在彗星带的探讨”. 皇家天文学会月报. 192 (3): 481–491. 文献标识码:1980MNRAS.192..481F. DOI:10.1093/mnras/192.3.481.
\item 邓肯；奎因；特雷梅恩（1988）。“短周期彗星的起源”. 天体物理学报. 328：L69。文献标识码:1988ApJ...328L..69D. doi:10.1086/185162.
\item 马斯登，B.S.；朱维特，D.；马斯登，B.G.（1993）。“1993 FW”。国际天文学联合会通告（5730）。小行星中心：1。文献标识码:1993IAUC.5730....1L.
\item 戴奇斯，普雷斯顿（2018年12月14日）。“关于柯伊伯带的10个须知”。美国国家航空航天局太阳系探索于2019年1月10日存档自原始网页。检索日期：2019年12月1日
\item “柯伊伯带20岁”。天体生物学杂志。2012年9月1日。于2020年10月30日从原文存档。检索日期：2019年12月1日
\item 沃森，保罗（2019年1月1日）。“成功穿越冥王星之外的区域，NASA探测器开始传回柯伊伯带天体的图像”。科学。美国科学促进会。于2022年10月8日存档自原始网页。检索日期：2019年12月1日
\item 克莱德·汤博，《最后的发言》，致编辑的信件，《天空与望远镜》, 1994年12月，第8页
\item “‘柯伊伯带’这一术语有何不妥之处？”.国际彗星季刊于2019年10月8日存档自原始网页。检索日期：2021年12月19日
\item M. C. 德·桑克蒂斯；M. T. 卡普里亚与A. 科拉迪尼（2001）。“柯伊伯带天体的热演化与分异”. 《天文学杂志》. 121 (5): 2792–2799. 文献标识码:2001AJ....121.2792D. DOI:10.1086/320385.
\item “探索太阳系的边缘”。美国科学家网。2003年。于2009年3月15日从原文存档。检索日期：2007年6月23日
\item 迈克尔·E·布朗；玛格丽特·潘（2004）。“柯伊伯带的平面” （PDF）. 《天文学杂志》. 127 (4): 2418–2423. 文献标识码:2004AJ....127.2418B. DOI:10.1086/382515. S2CID 10263724。于2020年4月12日从原始内容存档 （PDF）。
\item 佩蒂，让-马克；莫尔比代利，亚历山德罗；瓦尔塞基，乔瓦尼·B.（1998）。“大质量散射行星体与小天体带的激发”（PDF）.《伊卡洛斯》141（2）：367。文献标识码:1999Icar..141..367PDOI10.1006/icar.1999.6166. 于2007年8月9日从原始网页存档。检索日期：2007年6月23日
\item 卢尼恩，乔纳森·I.（2003）。“柯伊伯带”（PDF）。于2007年8月9日从原文存档。检索日期：2007年6月23日。
\item 朱维特，D.（2000年2月）。“经典柯伊伯带天体（CKBOs）”。于2007年6月9日从原始网页存档。检索日期：2007年6月23日。”
\item 默尔丁，P.（2000）。“柯伊伯带天体”。天文学与天体物理学百科全书. Bibcode:2000eaa..bookE5403.. doi:10.1888/0333750888/5403. ISBN 978-0-333-75088-9.
\item 埃利奥特，J. L.；等（2005）。“深黄道巡天：柯伊伯带天体与半人马星的搜寻。II. 动力学分类、柯伊伯带平面及核心群体”（PDF）。《天文学杂志》129（2）：1117–1162。文献编码:2005AJ....129.1117EDOI10.1086/427395存档于原始网页2013年7月21日检索日期2012年8月18日
\item “天体命名：小行星”。国际天文学联合会于2008年12月16日存档自原始网页。检索日期：2008年11月17日
\item 佩蒂，J.-M.；格拉德曼，B.；卡韦拉尔斯，J.J.；琼斯，R.L.；帕克，J.（2011）。“经典柯伊伯带内核的现实与起源”（PDF）。欧洲行星科学大会—美国天文学会行星科学分会联合会议（2011年10月2日至7日）。文献标识码:2011epsc.conf..722P已于2016年3月4日存档自原始网页。检索于2016年5月4日
\item 莱维森，哈罗德·F.；莫尔比戴利，亚历山德罗（2003）。“海王星迁移期间天体向外迁移导致柯伊伯带的形成”。自然. 426 (6965): 419–421. 文献标识码:2003Natur.426..419L. DOI:10.1038/nature02120. PMID 14647375. S2CID 4395099.
\item 斯蒂芬斯，丹尼斯·C.；诺尔，基思·S.（2006）。“利用NICMOS探测六对海王星外双星：冷经典盘中双星比例极高”。天文学杂志. 130 (2): 1142–1148. arXiv:astro-ph/0510130. 文献编码:2006AJ....131.1142S. DOI:10.1086/498715. S2CID 204935715.
\item 弗雷泽，韦斯利·C.；布朗，迈克尔·E.；莫尔比迪利，亚历山德罗；帕克，亚历克斯；巴蒂金，康斯坦丁（2014）。“柯伊伯带天体的绝对星等分布”。天体物理学报. 782 (2): 100. arXiv:1401.2157. 文献编码:2014ApJ...782..100F. doi:10.1088/0004-637X/782/2/100. S2CID 2410254.
\item 莱维森，哈罗德·F.；斯特恩，S·艾伦（2001）。“主柯伊伯带倾角分布的尺寸依赖性研究”。天文学杂志. 121 (3): 1730–1735. arXiv:astro-ph/0011325. Bibcode:2001AJ....121.1730L. doi:10.1086/319420. S2CID 14671420.
\item 莫尔比戴利，亚历山德罗（2005）。“彗星及其储库的起源与动力学演化”。arXiv:astro-ph/0512256.
\item 帕克，亚历克斯·H.；卡韦拉尔斯，J.J.；佩蒂，让-马克；琼斯，琳恩；格拉德曼，布雷特；帕克，乔尔（2011）。“七对超宽跨海王星双星的特征描述”。天体物理学杂志. 743 (1): 159. arXiv:1108.2505. Bibcode:2011AJ....141..159N. doi:10.1088/0004-6256/141/5/159. S2CID 54187134.
\item 利维森，哈罗德·F.；莫尔比戴利，亚历山德罗；范莱霍文，克里斯塔；戈梅斯，R.（2008）。“天王星和海王星轨道动力学不稳定期间柯伊伯带结构的起源”。国际行星科学杂志. 196 (1): 258–273. arXiv:0712.0553. 文献标识码:2008Icar..196..258L. DOI:10.1016/j.icarus.2007.11.035. S2CID 7035885.
\item “跨海王星天体列表”。小行星中心于2010年10月27日从原始网页存档。检索日期：2007年6月23日
\item 张；乔丹，A. B.；米利斯，R. L.；布伊，M. W.；瓦瑟曼，L. H.；埃利奥特，J. L. 等（2003）。“柯伊伯带中的共振占据：5:2共振与特洛伊共振的案例研究”。天文学杂志. 126 (1): 430–443. arXiv:astro-ph/0301458. 文献标识码:2003AJ....126..430C. DOI:10.1086/375207. S2CID 54079935.
\item Wm. 罗伯特·约翰斯顿（2007）。“海王星外天体”。于2019年10月19日存档自原始网页。检索日期：2007年6月23日
\item E.I. 蔡与M.E. 布朗（1999）。“凯克笔尖束巡天观测暗弱柯伊伯带天体”（PDF）。《天文期刊》118（3）：1411。arXiv：astro-ph/9905292文献标识码1999AJ....118.1411CDOI10.1086/301005S2CID8915427存档于2012年6月12日的原始网页。检索于2007年7月1日
\item 伯恩斯坦，G. M.；特林灵，D. E.；艾伦，R. L.；布朗，K. E.；霍尔曼，M.；马尔霍特拉，R.（2004）。“海王星外天体的大小分布”。天文期刊. 128 (3): 1364–1390. arXiv:astro-ph/0308467. Bibcode:2004AJ....128.1364B. doi:10.1086/422919. S2CID 13268096.
\item 迈克尔·布鲁克斯（2005）。“13件说不通的事”。NewScientistSpace.com于2018年10月12日存档自原始网页。检索日期：2018年10月12日
\item 戈弗特·席林（2008）。“行星X之谜”。《新科学家》于2015年4月20日存档自原始网页。检索日期：2008年2月8日
\item C. 德拉富恩特·马尔科斯 & R. 德拉富恩特·马尔科斯（2024年1月）。“越过外缘，迈向未知：柯伊伯悬崖之外的结构”。皇家天文学会月报快报527（1）（2023年9月20日出版）：L110 –L114arXiv:2309.03885文献标识码2024MNRAS.527L.110DDOI10.1093/mnrasl/slad132已于2023年10月28日存档自原始网页。已检索2023年9月28日
\item 钱伯斯，K. C.；等（2019年1月29日），泛星1号巡天, arXiv:1612.05560
\item 弗莱林，H. A.；等（2020年10月20日）。“泛星1号数据库与数据产品”. 天体物理学报增刊系列. 251 (1): 7. arXiv:1612.05243. 文献编码:2020ApJS..251....7F. doi:10.3847/1538-4365/abb82d. S2CID 119382318.
\item Pan-STARRS向全球发布最大规模数字巡天数据, 哈佛-史密森天体物理中心，2016年12月19日,已存档自原始页面于2022年10月21日获取2022年10月21日
\item “冥王星可能存在以氨为燃料的冰火山”。天文学杂志，2015年11月9日。于2016年3月4日从原始网页存档。”
\item 库齐，杰弗里·N.；霍根，罗伯特·C.；博特克，威廉·F.（2010）。“迈向小行星与柯伊伯带天体的初始质量函数”。国际行星科学杂志. 208 (2): 518–538. arXiv:1004.0270. 文献标识码:2010Icar..208..518C. DOI:10.1016/j.icarus.2010.03.005. S2CID 31124076.
\item 约翰森，A.；雅克，E.；库齐，J. N.；莫比德利，A.；古内尔，M.（2015）。“小行星形成的全新范式”。载于米歇尔，P.；德梅奥，F.；博特克，W.（编）。小行星IV。太空科学丛书。亚利桑那大学出版社。第471页。arXiv:1505.02941. 书目编码:2015aste.book..471J. DOI:10.2458/azu_uapress_9780816532131-ch025. ISBN 978-0-8165-3213-1. S2CID 118709894.
\item 内斯沃尼，大卫；尤丁，安德鲁·N.；理查森，德里克·C.（2010）。“由引力坍缩形成的柯伊伯带双星系统”。天文学杂志. 140 (3): 785–793. arXiv:1007.1465. Bibcode:2010AJ....140..785N. doi:10.1088/0004-6256/140/3/785. S2CID 118451279.
\item 汉森，K.（2005年6月7日）。“早期太阳系的轨道 shuffle”。地质时代于2007年9月27日从原始网页存档。检索日期：2007年8月26日
\item 齐加尼斯，K.；戈梅斯，R.；莫比戴利，亚历山德罗；莱维森，哈罗德·F.（2005）。“太阳系巨行星轨道结构的起源”。自然. 435 (7041): 459–461. 天文学编码:2005Natur.435..459T. DOI:10.1038/nature03539. PMID 15917800. S2CID 4430973.
\item 汤姆斯，E.W.；邓肯，M.J.；莱维森，哈罗德·F.（2002）。“木星与土星之间天王星和海王星的形成”。天文期刊. 123 (5): 2862–2883. arXiv:astro-ph/0111290. Bibcode:2002AJ....123.2862T. doi:10.1086/339975. S2CID 17510705.
\item 帕克，亚历克斯·H.；卡韦拉尔斯，J.J.（2010）。“海王星散射过程中双小行星的破坏”。天体物理学报快报.722(2)：L204 –L208arXiv：1009.3495文献编码2010ApJ...722L.204PDOI10.1088/2041-8205/722/2/L204S2CID119227937
\item 洛维特，R.（2010）。“柯伊伯带可能源于碰撞”。自然. doi:10.1038/news.2010.522.
\item 内斯沃尼，大卫；莫尔比戴利，亚历山德罗（2012）。“含四、五、六颗巨行星的早期太阳系不稳定性统计研究”。天文学杂志. 144 (4): 117. arXiv:1208.2957. Bibcode:2012AJ....144..117N. doi:10.1088/0004-6256/144/4/117. S2CID 117757768.
\item 内斯沃尼，大卫（2015）。“基于柯伊伯带天体倾角分布的海王星缓慢迁移证据”。天文期刊. 150 (3): 73. arXiv:1504.06021. 文献标识码:2015AJ....150...73N. doi:10.1088/0004-6256/150/3/73. S2CID 119185190.
\item 内斯沃尼，大卫（2015）。“跳跃海王星可解释柯伊伯带核心。”天文学杂志. 150 (3): 68. arXiv:1506.06019. Bibcode:2015AJ....150...68N. doi:10.1088/0004-6256/150/3/68. S2CID 117738539.
\item 弗雷泽，韦斯利；等（2017）。“所有诞生于柯伊伯带附近的原行星均以双星形式形成”。自然·天文学. 1 (4) 0088. arXiv:1705.00683. Bibcode:2017NatAs...1E..88F. doi:10.1038/s41550-017-0088. S2CID 118924314.
\item 沃尔夫，舒勒；道森，丽贝卡·I.；默里-克莱，露丝·A.（2012）。“ Neptune on Tiptoes：保持寒冷经典柯伊伯带的动力学历史”。天体物理学杂志. 746 (2): 171. arXiv:1112.1954. 文献标识码:2012ApJ...746..171W. DOI:10.1088/0004-637X/746/2/171. S2CID 119233820.
\item 莫尔比戴利，A.；加斯帕尔，H.S.；内斯沃尼，D.（2014）。“内侧冷柯伊伯带奇特偏心率分布的起源”。国际行星科学杂志. 232: 81–87. arXiv:1312.7536. 文献标识码:2014Icar..232...81M. DOI:10.1016/j.icarus.2013.12.023. S2CID 119185365.
\item 布朗，迈克尔·E.（2012）。“柯伊伯带天体的组成”。地球与行星科学年鉴. 40 (1): 467–494. arXiv:1112.2764. Bibcode:2012AREPS..40..467B. doi:10.1146/annurev-earth-042711-105352. S2CID 14936224.
\item 大卫·C·朱伊特与简·卢（2004）。“柯伊伯带天体（50000）夸瓦尔上的结晶水冰”（PDF）。自然432（7018）：731–3。文献标识码：2004Natur.432..731JDOI10.1038/nature03111PMID15592406S2CID4334385。于2007年6月21日从原始网页存档。检索日期：2007年6月21日
\item “在太阳系外围发现被逐出的小行星——欧洲南方天文台望远镜首次确认在柯伊伯带发现富含碳的小行星”.www.eso.org于2019年5月31日存档自原始页面。检索日期：2018年5月12日
\item 戴夫·朱维特（2004）。“柯伊伯带天体的表面”。夏威夷大学。于2007年6月9日从原始网页存档。检索日期：2007年6月21日
\item 朱维特，大卫；刘，简（1998）。“柯伊伯带的光学-红外光谱多样性” （PDF）. 天文学杂志. 115 (4): 1667–1670. 文献标识码:1998AJ....115.1667J. DOI:10.1086/300299. S2CID 122564418。于2020年4月12日从原始内容存档 （PDF）。
\item 朱维特，大卫·C.；刘，简·X.（2001）。“柯伊伯带天体的颜色与光谱”。天文期刊. 122 (4): 2099–2114. arXiv:astro-ph/0107277. 文献标识码:2001AJ....122.2099J. doi:10.1086/323304. S2CID 35561353.
\item 布朗，R. H.；克鲁克香克，D. P.；彭德尔顿，Y.；维德，G. J.（1997）。“柯伊伯带天体1993SC的表面组成”。科学. 276 (5314): 937–9. 文献标识码:1997Sci...276..937B. DOI:10.1126/science.276.5314.937. PMID 9163038. S2CID 45185392.
\item 黄，伊恩；布朗，迈克尔·E.（2017）。“柯伊伯带小型天体的双峰颜色分布”. 天文期刊. 153 (4): 145. arXiv:1702.02615. 文献编码:2017AJ....153..145W. doi:10.3847/1538-3881/aa60c3. S2CID 30811674.
\item 布朗，迈克尔·E.；布莱克，杰弗里·A.；凯斯勒，雅克琳·E.（2000）。“明亮柯伊伯带天体2000 EB173的近红外光谱”。天体物理学报. 543(2)：L163。文献标识码:2000ApJ...543L.163B. CiteSeerX 10.1.1.491.4308. doi:10.1086/317277. S2CID 122764754.
\item 利坎德罗；奥利瓦；迪马蒂诺（2001）。“NICS-TNG对跨海王星天体2000 EB173和2000 WR106的红外光谱观测”。天文学与天体物理学. 373(3)：L29。arXiv:astro-ph/0105434. 文献编码:2001A&A...373L..29L. DOI:10.1051/0004-6361:20010758. S2CID 15690206.
\item 格拉德曼，布雷特；等（2001年8月）。“柯伊伯带的结构”。天文期刊. 122 (2): 1051–1066. 天文学文献标识符:2001AJ....122.1051G. DOI:10.1086/322080. S2CID 54756972.
\item 皮捷娃，E. V.；皮捷夫，N. P.（2018年10月30日）。“基于行星与航天器运动的主小行星带和柯伊伯带质量”。天文学快报. 44 (89): 554–566. arXiv:1811.05191. Bibcode:2018AstL...44..554P. doi:10.1134/S1063773718090050. S2CID 119404378.
\item 内斯沃尼，大卫；沃克鲁赫利茨基，大卫；博特克，威廉·F.；诺尔，凯思；莱维森，哈罗德·F.（2011）。“观测到的双星比例限制了柯伊伯带碰撞磨蚀的范围”。天文学杂志. 141 (5): 159. arXiv:1102.5706. 文献标识码:2011AJ....141..159N. DOI:10.1088/0004-6256/141/5/159. S2CID 54187134.
\item 莫尔比戴利，亚历山德罗；内斯沃尼，大卫（2020）。“柯伊伯带：形成与演化”。海王星外太阳系。第25–59页。arXiv:1904.02980。doi10.1016/B978-0-12-816490-7.00002-3ISBN9780128164907S2CID102351398
\item 香克曼，C.；卡韦拉尔斯，J. J.；格拉德曼，B. J.；亚历山德森，M.；凯布，N.；佩蒂，J.-M.；班尼斯特，M. T.；陈毅廷；格温，S.；雅库比克，M.；沃尔克，K.（2016）。“OSSOS. II. 夸奥尔带散射族群绝对星等分布中的显著跃迁”. 天文学杂志. 150 (2): 31. arXiv:1511.02896. Bibcode:2016AJ....151...31S. doi:10.3847/0004-6256/151/2/31. S2CID 55213074.
\item 亚历山德森，迈克；格拉德曼，布雷特；卡韦拉尔斯，J.J.；佩蒂，让-马克；格温，斯蒂芬；香克曼，科克（2014）。“一项精心刻画并持续追踪的海王星外天体巡天调查：冥王星族小天体的大小分布及海王星特洛伊小行星的数量”. 天文学杂志. 152 (5): 111. arXiv:1411.7953. doi:10.3847/0004-6256/152/5/111. S2CID 119108385.
\item “哈勃发现有史以来最小的柯伊伯带天体”。哈勃网站，2009年12月。已归档自原始网页于2021年1月25日存档。检索日期：2015年6月29日。
\item 施利希廷，H. E.；奥菲克，E. O.；文茨，M.；萨里，R.；加尔-亚姆，A.；利维奥，M. 等人（2009年12月）。“基于档案数据中的恒星掩星现象观测到的一颗直径不足一公里的柯伊伯带天体”。自然. 462 (7275): 895–897. arXiv:0912.2996. 文献标识码:2009Natur.462..895S. DOI:10.1038/nature08608. PMID 20016596. S2CID 205219186.
\item 施利希廷，H. E.；奥菲克，E. O.；文茨，M.；萨里，R.；加尔-亚姆，A.；利维奥，M. 等人（2012年12月）。“利用恒星掩食测量亚千米级柯伊伯带天体的丰度”. 《天体物理学杂志》. 761 (2): 10. arXiv:1210.8155. 文献标识码:2012ApJ...761..150S. DOI:10.1088/0004-637X/761/2/150. S2CID 31856299. 150.
\item 多纳尔，亚历克斯；霍拉尼，米哈伊；巴根纳尔，弗兰；布兰特，蓬图斯；格朗迪，威尔；利塞，凯里；帕克，乔尔；波佩，安德鲁·R.；辛格，凯尔西·N.；斯特恩，S·艾伦；韦伯西瑟，安妮（2024年2月1日）。“新视野号维内西亚·伯尼学生尘埃计数器观测到接近60天文单位时的通量高于预期”. 天体物理学报快报. 961(2): L38.arXiv:2401.01230. Bibcode:2024ApJ...961L..38D. doi:10.3847/2041-8213/ad18b0.
\item “半人马星与散盘天体列表”。国际天文学联合会：小行星中心存档于2017年6月29日自原始网页。检索日期：2010年10月27日
\item 大卫·朱维特（2005）。“1000公里尺度的柯伊伯带天体”。夏威夷大学于2017年7月2日存档自原始网页。检索日期：2006年7月16日
\item 克雷格·B·阿格诺尔与道格拉斯·P·汉密尔顿（2006）。“海王星在双行星引力遭遇中捕获其卫星海卫一”（PDF）。自然441（7090）：192–194。天文学编码：2006Natur.441..192ADOI10.1038/nature04792PMID16688170S2CID4420518。于2007年6月21日从原始网页存档。检索日期：2007年10月29日
\item 昂克纳兹，泰蕾丝；卡伦巴赫，R.；欧文，T.；索坦，C.（2004）。海卫一、冥王星、半人马星及海王星外天体。施普林格出版社。ISBN978-1-4020-3362-9。检索日期：2007年6月23日。
\item 来源：喷气推进实验室小天体数据库
\item 迈克·布朗（2007）。“厄里斯的卫星迪斯诺米亚”。加州理工学院于2012年7月17日存档自原始网页。检索日期：2007年6月14日
\item “决议B5和B6”（PDF）。国际天文学联合会。2006年。于2009年6月20日从原始网页存档。检索日期：2011年9月2日。
\item 格兰迪，W.M.；诺尔，K.S.；布伊，M.W.；贝内基，S.D.；拉戈齐尼，D.；罗，H.G.（2019年12月）。“跨海王星双星Gǃkúnǁʼhòmdímà（(229762) 2007 UK)的相互轨道、质量和密度”（PDF）。《土星》334：30–38。文献编码2019Icar..334...30GDOI10.1016/j.icarus.2018.12.037S2CID126574999已归档于2019年4月7日自原始网页存档。
\item 迈克·布朗“外太阳系有多少颗矮行星？”存档于2011年10月18日，互联网档案馆访问日期：2013年11月15日
\item 坦克雷迪，G.；法夫尔，S. A.（2008）。“太阳系中的矮行星有哪些？”伊卡洛斯. 195 (2): 851–862. Bibcode:2008Icar..195..851T. doi:10.1016/j.icarus.2007.12.020.
\item 布朗，M. E.；范达姆，M. A.；布歇，A. H.；勒米尼耶，D.；坎贝尔，R. D.；陈，J. C. Y.；康拉德，A.；哈特曼，S. K.；约翰逊，E. M.；拉丰，R. E.；拉比诺维茨，D. L. 拉比诺维茨；斯通斯基，P. J. Jr.；萨默斯，D. M.；特鲁希略，C. A.；维齐诺维奇，P.L.（2006）。“柯伊伯带最大天体的卫星”（PDF）。天体物理学杂志639（1）：L43－L46arXiv：astro-ph/0510029文献标识码2006ApJ...639L..43BDOI10.1086/501524S2CID2578831于2018年9月28日存档自原始网页。检索于2011年10月19日
\item 阿格诺尔，C.B.；汉密尔顿，D.P.（2006）。“海王星在双行星引力遭遇中捕获其卫星海卫一”（PDF）。自然441（7090）：192–4。天文学编码：2006Natur.441..192ADOI10.1038/nature04792PMID16688170S2CID4420518已归档于2013年11月3日自原始网页存档。检索于2010年7月9日
\item “美国宇航局新视野号团队发布首次柯伊伯带飞越科学成果”。美国宇航局。2019年5月16日。已归档自原始网页于2019年12月16日存档。检索日期：2019年5月16日。
\item “新视野计划：新地平线科学目标”。美国国家航空航天局——新视野计划。2015年4月15日存档自原文。检索日期：2015年4月15日
\item “美国宇航局哈勃望远镜为新视野号冥王星任务发现潜在的柯伊伯带目标”。新闻稿。约翰霍普金斯应用物理实验室。2014年10月15日。于2014年10月16日从原文存档。检索于2014年10月16日
\item 布伊，马克（2014年10月15日）。“新视野号哈勃望远镜柯伊伯带天体搜寻结果：状态报告”（PDF）。太空望远镜科学研究所。第23页。于2015年7月27日从原文存档。检索日期：2015年8月29日
\item 拉克达瓦拉，艾米丽（2014年10月15日）“终于！新视野号有了第二个目标！”行星协会博客行星协会已存档自原始网页于2014年10月15日检索日期2014年10月15日
\item “哈勃望远镜将继续全面搜寻新视野目标”。哈勃网站新闻稿太空望远镜科学研究所，2014年7月1日。已存档于2015年5月12日自原始页面。检索日期：2014年10月15日
\item 斯特罗姆伯格，约瑟夫（2015年4月14日）。“美国宇航局的新视野号探测器正在探访冥王星——并刚刚传回了首批彩色照片”。Vox于2020年4月6日从原始页面存档。检索日期：2015年4月14日
\item 科里·S·鲍威尔（2015年3月29日）：“艾伦·斯特恩谈冥王星的奇观、新视野号的失联双子星，以及那所谓的‘矮行星’问题”。发现。于2019年11月16日存档自原文。检索日期：2015年8月29日
\item 波特，S. B.；帕克，A. H.；布伊，M.；斯宾塞，J.；韦弗，H.；斯特恩，S. A.；贝内基，S.；赞加里，A. M.；韦伯西尔，A.；吉温，S.；佩蒂，J. -M.；斯特纳，R.；博恩坎普，D.；诺尔，K.；卡维拉尔斯，J. J.；托伦，D.；辛格，K. N.；肖沃尔特，M.；富恩特斯，C.；伯恩斯坦，G.；贝尔顿，M.（2015）。“潜在新视野号柯伊伯带天体探测目标的轨道与可及性” （PDF）. 美国宇航研究协会-休斯顿 (1832): 1301. 文献标识码:2015LPI....46.1301P。于2016年3月3日从原文 （PDF）存档。
\item 麦金农，米卡（2015年8月28日）。“新视野号锁定下一个目标：让我们探索柯伊伯带吧！”。地球与空间于2015年12月31日存档自原始网页。
\item 德韦恩·布朗／劳里·坎蒂略（2016年7月1日）。“新视野号获准延长任务至柯伊伯带，黎明号将继续驻留谷神星”。美国国家航空航天局。于2016年8月20日从原始页面存档。检索日期：2017年5月15日。
\item 新视野号探测到一颗距离不远的流浪柯伊伯带天体 已归档2021年11月26日于互联网档案馆spacedaily.com，马里兰州劳雷尔（SPX）。2015年12月7日。
\item 科鲁姆，乔纳森（2019年2月10日）。“新视野号瞥见了终极塞勒斯的扁平形状——美国国家航空航天局的新视野号探测器飞越了有史以来最遥远的天体：一个来自太阳系早期的微小天体，名为2014 MU69，昵称‘终极塞勒斯’。”– 互动.纽约时报于2021年12月24日从原始版本存档。检索日期：2019年2月11日
\item 霍尔，劳拉（2017年4月5日）。“支持聚变的冥王星轨道器与着陆器”。美国国家航空航天局于2017年4月21日从原始页面存档。检索日期：2018年7月13日
\item “全球航天公司拟向美国国家航空航天局展示冥王星着陆器概念”。EurekAlert!于2019年1月21日存档自原始网页。检索日期：2018年7月13日
\item 庞西，乔尔；丰德卡巴·拜格，若尔迪；费雷辛，弗雷德；马蒂诺，文森特（2011年3月1日）。“对海梅亚系统中轨道器的初步评估：行星轨道器能多快抵达如此遥远的目标？”航天学报. 68 (5–6): 622–628. 文献编码:2011AcAau..68..622P. DOI:10.1016/j.actaastro.2010.04.011. ISSN 0094-5765.
\item “哈乌梅亚：技术与原理”.www.centauri-dreams.org于2018年7月13日存档自原始网页。检索日期：2018年7月13日
\item 《新视野号飞向冥王星的震撼历程在新书中揭秘》。Space.com于2018年7月13日存档自原始页面。检索日期：2018年7月13日
\item TVIW（2017年11月4日）22. 人类迈向另一颗恒星的第一步：星际探测任务，于2021年10月30日存档自原文，检索于2018年7月24日
\item “三年一度的地球—太阳峰会”。已归档自原始网页，日期为2020年8月3日。检索日期：2018年7月24日
\item 格利夫斯，阿什利；艾伦，兰德尔；图皮斯，亚当；奎格利，约翰；穆恩，亚当；罗伊，埃里克；斯宾塞，大卫；尤斯特，尼古拉斯；莱恩，詹姆斯（2012年8月13日）。对跨海王星天体任务机遇的调查——第二部分，轨道捕获。美国航空航天学会/美国宇航学会天体力学专题会议，明尼苏达州明尼阿波利斯。弗吉尼亚州雷斯顿：美国航空航天学会。DOI:10.2514/6.2012-5066. ISBN 9781624101823. S2CID 118995590.
\item 卡拉斯，保罗；格雷厄姆，詹姆斯·R.；克兰平，马克·C.；菲茨杰拉德，迈克尔·P.（2006）。“HD 53143和HD 139664周围尘埃盘的首批散射光图像”。天体物理学报. 637(1)：L57。arXiv:astro-ph/0601488. 文献标识码:2006ApJ...637L..57K. doi:10.1086/500305. S2CID 18293244.
\item 特林宁，D. E.；布莱登，G.；贝希曼，C. A.；里克，G. H.；苏，K. Y. L.；斯坦斯伯里，J. A.；布莱洛克，M.；斯塔彭费尔特，K. R.；比曼，J. W.；哈勒，E. E.（2008年2月）。“类太阳恒星周围的尘埃盘”。天体物理学杂志. 674 (2): 1086–1105. arXiv:0710.5498. Bibcode:2008ApJ...674.1086T. doi:10.1086/525514. S2CID 54940779.
\item “两颗邻近恒星周围的尘埃行星盘酷似我们的柯伊伯带”，2006年。于2016年7月9日存档自原始网页。检索日期：2007年7月1日。
\item 库赫纳，M. J.；斯塔克，C. C.（2010）。“柯伊伯带尘埃云的碰撞清理模型”。天文学杂志. 140 (4): 1007–1019. arXiv:1008.0904. Bibcode:2010AJ....140.1007K. doi:10.1088/0004-6256/140/4/1007. S2CID 119208483.
\end{enumerate}
\subsection{外部链接}
\begin{itemize}
\item 戴夫·朱维特在加州大学洛杉矶分校的主页
\item 该带的名称
\item 按族分类的短周期彗星列表 已归档2019年6月24日于互联网档案馆
\item 柯伊伯带简介由美国国家航空航天局太阳系探索提供
\item 柯伊伯带电子通讯
\item Wm. 罗伯特·约翰斯顿的外海天体页面
\item 小行星中心：外太阳系分布图, 用于展示柯伊伯空隙
\item 外海天体列表
\end{itemize}